\documentclass{article}

\PassOptionsToPackage{numbers,compress}{natbib}
\usepackage{neurips_2026}

\usepackage[utf8]{inputenc}
\usepackage[T1]{fontenc}
\usepackage{hyperref}
\usepackage{url}
\usepackage{booktabs}
\usepackage{amsfonts}
\usepackage{amsmath}
\usepackage{microtype}
\usepackage{graphicx}

\newcommand{\auc}{\mathrm{AUC}}

\title{Trial-Level Cross-Validation Can Invert\\
Within- Versus Cross-Subject Comparisons}

\author{%
  Anonymous Author(s)\\
  Affiliation\\
  \texttt{email}\\
}

\begin{document}
\maketitle

\begin{abstract}
A common design in neural decoding compares a model trained on a subject's own data against one
trained on other subjects. The within-subject arm is usually estimated by cross-validating over
trials. We show that when a recording session contains structure that happens to correlate with the
label, this estimate is inflated -- and, because the cross-subject arm shares no session with its
test data and cannot be inflated the same way, the bias falls entirely on one side of the
comparison. In simulation with known ground truth the bias reverses the sign of the conclusion: at a
moderate nuisance level a trial-level split reports the within-subject decoder beating cross-subject
by $+0.026$ $\auc$ when it in fact loses by $0.008$. A label-\emph{independent} session offset
produces no such bias, so the mechanism is specifically session-level correlation with the label, not
session identity. Nonlinear models exploit it $1.7\times$ more than linear ones, while $\ell_2$
strength has no effect. On a widefield calcium cohort (25 mice, DANDI:001425) the same substitution
removes $30$--$73\%$ of a published-style asymmetry and eliminates one of four effects entirely,
turning $+0.014$ $\auc$ at $p = 0.001$ into $+0.003$ at $p = 0.37$. On a second, independent
widefield cohort we find no such inflation, so the bias is a conditional hazard rather than a
universal one. We give the condition, the diagnostic, and the fix.
\end{abstract}

\section{The problem}
Many decoding studies ask whether a model trained on other individuals can substitute for one
trained on the subject, and answer it by comparing two accuracies: a within-subject estimate and a
cross-subject estimate. The comparison looks symmetric. It is not.

Trials recorded in one session are not independent. They share illumination and hemodynamic state,
slow arousal, electrode or lens drift, and simple temporal autocorrelation. A within-subject decoder
cross-validated over trials can therefore be trained on part of a session and tested on the rest of
the same session. A cross-subject decoder is trained on other individuals entirely and never shares
a session with its test data. Whatever advantage the shared session confers accrues to one arm only,
and therefore lands directly on the difference the study reports.

This note establishes three things: the precise structure that causes the bias, how large it gets,
and when it does \emph{not} occur.

\section{The mechanism, with ground truth}
\label{sec:sim}
We simulate $12$ subjects $\times$ $5$ sessions $\times$ $100$ trials in $44$ dimensions:
\begin{equation}
x = \delta\, y\, w \;+\; \gamma\, y\, z_{s,k} \;+\; u_s \;+\; v_{s,k} \;+\; \varepsilon,
\qquad \varepsilon \sim \mathcal{N}(0, I).
\end{equation}
Here $w$ is a decodable direction shared by all subjects, $u_s$ a subject offset, $v_{s,k}$ a
per-session offset of scale $\sigma$, and $z_{s,k}$ a direction drawn fresh for each session that
correlates with the label \emph{within that session only}. Labels are balanced Bernoulli, so
$v_{s,k}$ carries no label information by construction and any accuracy attributable to it is
leakage. We compare three estimators: within-subject with a trial-level 5-fold, within-subject with
leave-one-session-out, and cross-subject (train on the other 11 subjects).

\textbf{Session identity alone leaks nothing.} Sweeping $\sigma$ from $0$ to $2$ with $\gamma = 0$
produces inflation between $-0.014$ and $+0.012$ $\auc$, i.e.\ nothing systematic. This is worth
stating because it is the intuitive story and it is wrong: a constant offset shifts a session's
cluster without moving the optimal boundary, so a linear decoder gains nothing from having seen the
session before.

\textbf{Session-level label correlation does.} Holding $\sigma = 0.5$ and sweeping $\gamma$
(Table~\ref{tab:sim}), inflation is absent at $\gamma \leq 0.5$ and then grows sharply. The
consequence is not merely that a number is too large. At $\gamma = 1$ the trial-level split reports
the within-subject decoder ahead of cross-subject by $+0.026$ $\auc$; leave-one-session-out shows it
is behind by $0.008$. The bias reverses the sign of the reported comparison. At $\gamma = 2$ it
manufactures $+0.115$ from a true $-0.020$.

\begin{table}[t]\centering
\caption{Simulation with known ground truth ($\sigma = 0.5$, $T = 100$, $\ell_2$ logistic).
\emph{Inflation} is within-trial minus within-session. The last two columns are the within-minus-cross
gap a study would report under each protocol; they differ in sign from $\gamma = 1$ onward.}
\label{tab:sim}
\begin{tabular}{lccccc}
\toprule
$\gamma$ & within (trial) & within (session) & cross & inflation & reported gap: trial / session \\
\midrule
0.00 & 0.596 & 0.610 & 0.621 & $-0.014$ & $-0.025$ / $-0.011$ \\
0.25 & 0.601 & 0.613 & 0.620 & $-0.012$ & $-0.019$ / $-0.007$ \\
0.50 & 0.611 & 0.613 & 0.619 & $-0.003$ & $-0.008$ / $-0.005$ \\
1.00 & 0.642 & 0.608 & 0.616 & $\mathbf{+0.034}$ & $\mathbf{+0.026}$ / $\mathbf{-0.008}$ \\
2.00 & 0.723 & 0.589 & 0.609 & $\mathbf{+0.135}$ & $\mathbf{+0.115}$ / $\mathbf{-0.020}$ \\
\bottomrule
\end{tabular}
\end{table}

\textbf{Model class matters; regularization does not.} At $\gamma = 1$, inflation is $+0.0338$,
$+0.0337$ and $+0.0338$ for $\ell_2$ logistic regression at $C = 0.01$, $0.5$ and $100$ -- identical
to four decimals across four orders of magnitude of penalty. A one-hidden-layer network on the same
data inflates by $+0.0564$, $1.7\times$ more. The quantity that matters is the flexibility of the
model class, not how hard it is penalized within a class. A capacity comparison run on a trial-level
split therefore confounds capacity to fit the signal with capacity to exploit the leak.

\section{A real instance}
\label{sec:braidyn}
BraiDyn-BC \citep{kondo2025} (DANDI:001425) provides 25 mice $\times$ $\sim$15 daily sessions of
cortex-wide widefield $\Delta F/F$ in 44 Allen-atlas parcels. We decode four task events and measure
the personal-minus-pooled aging asymmetry -- exactly the within- versus cross-subject shape described
above -- under both protocols on identical features.

Substituting leave-one-session-out for the trial-level split lowers the within-subject arm by
$+0.013$, $+0.010$, $+0.009$ and $+0.000$ $\auc$ across the four events, and leaves the cross-subject
arm untouched, as predicted. The effect on the reported comparison (Table~\ref{tab:braidyn}) is
substantial: $30$--$73\%$ of the asymmetry disappears, and for lever pull the entire effect does,
turning $+0.014$ at $p = 0.001$ into $+0.003$ at $p = 0.37$. The remaining events survive at reduced
magnitude.

The capacity prediction of Section~\ref{sec:sim} also holds here. Holding out sessions removes
$43$--$73\%$ of the CNN's asymmetry against $30$--$48\%$ of the linear model's on the same events, so
the higher-capacity model was the larger beneficiary of the leak -- and a naive reading of the
trial-split numbers would have concluded that the effect \emph{grows} with capacity, which it does
not.

\begin{table}[t]\centering
\caption{BraiDyn-BC, $n = 25$ mice. Same features, same events, same statistic; only the split
differs.}
\label{tab:braidyn}
\begin{tabular}{lcc}
\toprule
Event & trial-level split & session-held-out \\
\midrule
Lever-pull & $+0.014$ ($p = 0.001$) & $+0.003$ ($p = 0.37$) \\
Lick       & $+0.028$ ($p = 0.0002$) & $+0.020$ ($p = 0.002$) \\
Reward     & $+0.032$ ($p = 0.0001$) & $+0.029$ ($p = 0.0005$) \\
Tone       & $+0.022$ ($p = 0.005$) & $+0.015$ ($p = 0.08$) \\
\midrule
Pooled     & $+0.024$ & $+0.017$ \\
\bottomrule
\end{tabular}
\end{table}

\section{A cohort where it does not occur}
\label{sec:cardin}
We ran the same substitution on an independent widefield cohort \citep{benisty2024}: a different
laboratory, a different parcellation (23 CCFv3 regions), spontaneous rather than cued behavior, and
locomotion and pupil-derived arousal as targets. Each animal is one continuous recording, so the
analogue of a session-held-out split is a contiguous-block split.

We found no inflation. For locomotion the within-subject arm changed by $-0.005$ $\auc$ and for
arousal by $-0.046$; in both cases the shuffled estimate was \emph{lower} than the blocked one, and
for arousal this held in $6$ of $6$ animals. The within-minus-cross gap was accordingly not
overstated in this cohort.

We report this because it bounds the claim. The bias is not an automatic consequence of splitting
trials; it requires the session to contain label-correlated structure for the decoder to absorb.
Section~\ref{sec:sim} makes that condition explicit, and a cohort can fail to meet it.

\section{Recommendation}
Hold out whole recording sessions when estimating the within-subject arm of a within- versus
cross-subject comparison. The cost is a slightly noisier estimate; the benefit is that the two arms
are exposed to the same information.

Because the hazard is conditional, we also suggest reporting both numbers when the comparison is
load-bearing. The difference between them is itself the diagnostic: it is an estimate of how much
session-level label-correlated structure the data contain, costs one extra fit to obtain, and would
have flagged the effect in Section~\ref{sec:braidyn} immediately.

\textbf{Limitations.} We have one positive cohort and one negative one, so we cannot say how common
the hazard is in practice, only that it is real, that it can invert a conclusion, and that it is not
universal. Our account of why the second cohort is unaffected is post hoc: we predicted inflation
there and did not find it, and while the mechanism in Section~\ref{sec:sim} is consistent with the
null, we have not independently measured label-correlated session structure in either cohort and
shown it differs. A direct measurement of $\gamma$ in real data, and a survey of published
within-versus-cross designs, are the natural next steps.

{\small
\bibliographystyle{plainnat}
\begin{thebibliography}{9}
\bibitem[Benisty et al.(2024)]{benisty2024} H.~Benisty et al. Rapid fluctuations in functional
connectivity of cortical networks encode spontaneous behavior. \emph{Nature Neuroscience}, 27, 2024.
figshare 175317.
\bibitem[Kondo et al.(2025)]{kondo2025} A.~Kondo et al. BraiDyn-BC: a cortex-wide widefield calcium
imaging dataset during a cued lever-pull operant task. \emph{Scientific Data}, 12, 2025.
DANDI:001425.
\end{thebibliography}
}

\end{document}
